% Compile beside math-review.sty; replace all visible insertion text for publication.
% Rename and regroup functional headings around the actual mathematical thread.
% Shared files do not determine section hierarchy; no environment/count quota applies.
\documentclass[12pt,reqno]{amsart}
\usepackage{math-review}
\title[Part IV: Boundaries (standard)]{[Review topic]\\
Part IV: Applications, Boundaries and Problems (standard)}
\author{[Author name]}
\date{[Prepared date; literature cutoff date]}
\begin{document}
\begin{abstract}
\placeholder{Identify the applications, boundaries and core Problems, emphasizing the deeper explanations of progress, remaining ranges and mathematical relationships.}
\end{abstract}
\maketitle
\tableofcontents

\section{Introduction}\label{sec:introduction}
\placeholder{Explain the central question about the theory's applicability and remaining frontier. Present the larger mathematical themes linking its achievements, limits and unresolved goals.}
\placeholder{Explain how the applications in Section~\ref{sec:applications} illuminate the boundaries in Section~\ref{sec:boundaries}, and how these lead to the Problems in Section~\ref{sec:problems}. Identify what deeper analysis of uses and progress will add.}

\section{Applications that reveal the force of the results}\label{sec:applications}
\placeholder{Develop meaningful uses or constructions with their exact settings. Verify assumptions, derive conclusions and explain why the result is useful. Deepen core examples before adding variants.}

\subsection{A structural feature behind the application}\label{subsec:feature}
\placeholder{Explain what makes the application work and which part of the argument transfers. Retain this local development only when it supports the larger application question.}

\section{Established boundaries and unsettled extensions}\label{sec:boundaries}
\placeholder{Explain exact limitations, exceptional cases or known counterexamples where available. Compare changed hypotheses carefully. Separate failure of a method from impossibility and state the precise gap that motivates the next section.}

\section{Core Problems, progress and relationships}\label{sec:problems}
\placeholder{Group the agreed frontier around its mathematical targets. Explain the common issue before treating local Problems; do not substitute an exhaustive list for a theory of their relationships.}
\begin{problem}[A core unresolved target]\label{prob:core}
\placeholder{Give the exact objects, domain, hypotheses, quantifiers and goal from the source formulation. Retain decisive conditions and distinguish a source Problem from a review-proposed direction.}
\end{problem}
\placeholder{Give the source and cutoff status after Problem~\ref{prob:core}. Develop its importance, verified advances, exact solved parameter ranges and the remaining scope. Explain relevant difficulties or reductions with evidence; do not invent a strategy or claim every unmentioned case is open.}

\subsection{What progress changes about the Problem}\label{subsec:progress}
\placeholder{Compare a significant verified partial result with the full target and explain what it resolves and leaves. Justify relationships with other core Problems, and distinguish sourced deductions from speculation. Use the same conditions and status wherever this Problem is shared.}
\placeholder{A narrow settled topic needs an explicit, verified settlement account rather than invented open Problems. Discuss neighboring questions only with their different scope identified.}

\templatereferences
\end{document}
